\section{Method}

\subsection{Problem Setup}

Each observation consists of an RGB image, a visible partial point cloud extracted from a registered global scan, graph connectivity over that partial point cloud, and a small vector of 2D and 3D visibility statistics. The full scene is also represented as a coarsened global point-cloud graph. We study two tasks:
\begin{enumerate}[leftmargin=*]
  \item learning a shared latent space between image observations and partial geometry; and
  \item localizing latent observations against the larger scene graph.
\end{enumerate}

\subsection{Stage 1: Visibility-Aware Cross-Modal Autoencoding}

The image branch uses a convolutional encoder with token pooling, while the geometry branch uses an MLP projection followed by graph convolutions and global pooling. Both branches produce a latent vector of matching dimensionality. From these latents, the model reconstructs:
\begin{enumerate}[leftmargin=*]
  \item image from image latent,
  \item image from point-cloud latent,
  \item point cloud from point-cloud latent,
  \item point cloud from image latent.
\end{enumerate}

This setup encourages cross-modal alignment and supports a latent consistency loss between the image and geometry encoders.

\subsection{Stage 2: Random-Walk Observation Encoding}

To generate sequential supervision, the repository synthesizes random walks through manually defined navigable regions. Each waypoint yields a rendered RGB observation, a visible point-cloud slice, and camera-relative metadata such as location and view direction. The trained autoencoder is then used to encode these observations into latent trajectories.

\subsection{Stage 3: Localization on a Coarsened Global Scene Graph}

The full aligned scan is downsampled, converted into a radius graph, and partitioned into local clusters. A localization network then compares random-walk latents against these cluster embeddings. The current implementation predicts:
\begin{enumerate}[leftmargin=*]
  \item observation-to-cluster translation vectors, and
  \item relative transforms between nearby latent observations.
\end{enumerate}

This turns localization into graph-conditioned metric reasoning rather than direct image-only regression.

\subsection{Training Objectives}

The current losses in the code suggest the following objective:
\begin{align}
\mathcal{L} = &\ \lambda_{\text{img}} \left(\mathcal{L}_{\text{img}\rightarrow\text{img}} + \mathcal{L}_{\text{pcd}\rightarrow\text{img}}\right) \\
&+ \lambda_{\text{pcd}} \left(\mathcal{L}_{\text{pcd}\rightarrow\text{pcd}} + \mathcal{L}_{\text{img}\rightarrow\text{pcd}}\right) \\
&+ \lambda_{\text{align}} \mathcal{L}_{\text{latent}} + \lambda_{\text{loc}} \mathcal{L}_{\text{localization}}.
\end{align}

For geometry reconstruction, the code currently uses Chamfer distance plus a nearest-neighbor color term. For cross-modal alignment, it uses mean-squared error between latent vectors. For localization, it uses MSE on observation-to-cluster translations and pairwise relative transforms.

\subsection{Intended Paper Framing}

The strongest methodological framing is:
\begin{enumerate}[leftmargin=*]
  \item visibility-aware preprocessing creates structured geometric supervision,
  \item cross-modal autoencoding learns observation latents tied to geometry, and
  \item graph-based localization tests whether the latent space preserves global spatial structure.
\end{enumerate}
